
\documentclass[11pt]{article}

% Basic packages
\usepackage{fullpage}
\usepackage{amsmath,amssymb,amsthm}
\usepackage{hyperref}
\usepackage{mathrsfs}
\usepackage{enumitem}
\usepackage{algorithm}
\usepackage{algorithmic}
\usepackage{algpseudocode}
\usepackage{graphicx}
\usepackage{booktabs}
\usepackage{url}
\usepackage{xcolor}

% Theorem environments
\newtheorem{theorem}{Theorem}
\newtheorem{lemma}{Lemma}
\newtheorem{definition}{Definition}

%%%%%%%%%%%%%%%%%%%%%%%%%%%%%%%%%%%% Simple Notation
\newcommand{\Q}{{\mathcal{Q}}}
\newcommand{\K}{{\mathcal{K}}}
\newcommand{\V}{{\mathcal{V}}}
\newcommand{\D}{{\mathcal{D}}}
\newcommand{\Wq}{{W^{\Q}_i}}
\newcommand{\Wk}{{W^{\K}_i}}
\newcommand{\Wv}{{W^{\V}_i}}
\newcommand{\Wo}{{W^O}}
\newcommand{\h}{{\text{head}}}


%%%%%%%%%%%%%%%%%%%%%%%%%%%%%%%%%%%%% Comment
\newcommand{\krev}[1]{{\color{magenta}{#1}}} % 민혜
\newcommand{\brev}[1]{{\color{blue}{#1}}} % 병규
\begin{document}

\subsection{Notation of Transformer}
where $n$ is sequence length, $d_m$ is the embedding demension, and $h$ is the number of heads, $d_k$, $d_v$ are the hidden dimensions of the projection subspaces. For the rest of this paper, we will not differentiate
between $d_k$ and $d_v$ and just use $d$.
\begin{itemize}
    \item $\Q, \K, \V \in \mathbb{R}^{n\times d_m}$ : input embedding matrices
    \item $\Wq, \Wk \in \mathbb{R}^{d_m\times d_k}, \Wv \in \mathbb{R}^{d_m\times d_v}$: learned matrices
    \item \krev{$\Q\Wq = \mathbf{Q}, \K\Wk = \mathbf{K}, \V\Wv = \mathbf{V}$}
    \item \krev{$\Q\Wq^{'} = \mathbf{Q}^{'}, \K\Wk^{'} = \mathbf{K}^{'}, \V\Wv^{'} = \mathbf{V}^{'}$}
\end{itemize}

The Transformer is built upon the idea of Multi-Head Self-Attention (MHA), which allows the model to jointly attend to information at different positions from different representation subspaces. MHA is defined as:
\begin{equation} \label{multiattn}
\text{MultiHead}(\Q, \K, \V) = \text{Concat}(\h_1, \h_1, \cdots, \h_h)\Wo,
\end{equation}

Each head is defined as:
\begin{equation} \label{attn}
\text{head}_i = \text{Attention}(\Q\Wq,\K\Wk, \V\Wv) = \underbrace{\text{softmax}\left[\frac{\Q\Wq(\K\Wk)^\top}{\sqrt{d_k}}\right]}_P\V\Wv
\end{equation}
The self-attention defined in Equation~\ref{attn} refers to a context mapping matrix $P \in \mathbb{R}^{n\times n}$. The Transformer
uses $P$ to capture the input context for a given token, based on a combination of all tokens in the sequence. However, computing $P$ is expensive. It requires multiplying two $n \times d$ matrices, which is $O(n^2)$ in time and space complexity. This quadratic dependency on the sequence length has become a bottleneck for Transformers.

\subsection{Notation for Linformer}
Linformer adds two linear projection matrices $E_i,F_i \in \mathbb{R}^{n\times k}$ when computing key and value. They first project the original ($n\times d$)-dimensional key and value layers $\K\Wk$ and $\V\Wv$ into ($k \times d$)-dimensional projected key and value layers. Then they compute an ($n\times k$)-dimensional context mapping matrix $\bar{P}$ using scaled dot-product attention.
\begin{equation} \label{linattn}
\overline{\h_i} = \text{Attention}(\Q\Wq,E_i\K\Wk, F_i\V\Wv) = \underbrace{\text{softmax}\left[\frac{\Q\Wq(E_i\K\Wk)^\top}{\sqrt{d_k}}\right]}_{\bar{P}:n\times k}\underbrace{F_i\V\Wv}_{k\times d}
\end{equation}
Finally, Linformer compute context embeddings for each $\head_i$ using $\bar{P} \cdot (F_i\V\Wv)$. Note the above operations only require $O(nk)$ time and space complexity. Thus, if they can choose a very small projected dimension $k$, such that $k << n$, then we can significantly reduce the memory and space consumption. The following theorem states that, when $k=O(d/\epsilon^2)$ (independent of $n$), one can approximate $P\cdot \V\Wv$ using linear self-attention above Equation~\ref{linattn} with $\epsilon$ error.

\subsection{Notation for Performer}
Since Performer performs approximation including the softmax, Equation~\ref{attn} can be reformulated as follows:
\begin{equation} \label{perattn}
\begin{split}
\text{head}_i = \text{Attention}(\Q\Wq,\K\Wk, \V\Wv) = D^{-1}A\V\Wv, \\ A=\text{exp}(\Q\Wq(\K\Wk)^\top/\sqrt{d_k}), \\ D=\text{diag}(A1_n)
\end{split}
\end{equation}
Here $\text{exp}(\cdot)$ is applied elementwise, $1^n$ is the all-ones vector of length $n$, and $\text{diag}(\cdot)$ is a diagonal matrix with the input vector as the diagonal.

Performer works for attention blocks using matrices $A \in \mathbb{R}^{n\times n}$ of the form $A(i,j)= \textbf{K}(q^\top_i, k^\top_i)$, with $q_i/k_j$ standing for the $i^{th}/j^{th}$ query/key row-vector in $\Q\Wq/\K\Wk$ and kernel $\textbf{K} : \mathbb{R}^d \times \mathbb{R}^d \rightarrow \mathbb{R}_{+}$ defined for the (usually randomized) mapping: $\phi : \mathbb{R}^d \rightarrow \mathbb{R}^k_{+}$ (for some $k>0$) as:
\begin{equation} \label{kmapping}
\begin{split}
\textbf{K}(x,y)=\mathbb{E}[\phi(x)^\top \phi(y)]
\end{split}
\end{equation}
they call $\phi(u)$ a \textit{random feature map} for $u \in \mathbb{R}^d$. For $\Q\Wq^{'}, \K\Wk^{'} \in \mathbb{R}^{n\times k}$ with rows given as $\phi(q_i^\top)^\top$ and $\phi(k_i^\top)^\top$ respectively, Equation \label{perattn} leads directly to the efficient attention mechanism of the form:
\begin{equation} \label{perattnapprox}
\begin{split}
\widehat{\text{head}_i} = \text{Attention}(\Q\Wq,\K\Wk, \V\Wv) = \widehat{D^{-1}}(\Q\Wq^{'}((\K\Wk^{'})^\top \V\Wv)), \\ \widehat{D}=\text{diag}(\Q\Wq^{'}((\K\Wk^{'})^\top 1_n))
\end{split}
\end{equation}
Here $\widehat{\text{head}_i}$ stands for the approximate attention and brackets indicate the order of computations. It is easy to see that such a mechanism is characterized by space complexity $O(nk + nd + kd)$ and time complexity $O(nkd)$ as opposed to $O(n^2 + nd)$ and $O(n^2d)$ of the regular attention.


%mapping function에 대한 Notation
It turns out that by taking $\phi$ of the following form for functions $f_1,..., f_l : \mathbb{R} \rightarrow \mathbb{R}$, deterministic vectors $\omega_i$ or $\omega_1,...,\omega_m \overset{iid}{\sim} \mathcal{D}$ for some distribution $\mathcal{D} \in \mathcal{P}(\mathbb{R}^d)$:

\begin{equation} \label{mappingfunc}
\begin{split}
\phi(x) = h(x)\frac{1}{\sqrt{m}}(f_1(\omega^\top_1x),...,f_1(\omega^\top_mx),...,f_l(\omega^\top_1x),...,f_l(\omega^\top_mx))
\end{split}
\end{equation}
Performer can model most kernels used in practice. Furthermore, in most cases $\mathcal{D}$ is isotropic, usually Gaussian. Configurations: $h(x) = 1, l = 2, f_1 = sin, f_2 = cos$ correspond to shift-invariant kernels, in particular $\mathcal{D} = \mathcal{N}(0, \textbf{I}_d)$ leads to the Gaussian kernel  $\textbf{K}_{gauss}$(Rahimi & Recht, 2007). The softmax-kernel which defines regular attention matrix $A$ is given as:
\begin{equation} \label{smformal}
\begin{split}
\text{SM}(x,y) \overset{def}{=} \text{exp}(x^\top y)
\end{split}
\end{equation}
In the above, without loss of generality, they omit $\sqrt{d}$ renormalization since we can equivalently renormalize input keys and queries. Since: $\text{SM}(x,y) = \text{exp}(\frac{||x||^2}{2})K_{gauss}(x,y)\text{exp}(\frac{||y||^2}{2})$, based on what Performer have said, we obtain random feature map unbiased approximation of $\text{SM}(x,y)$ using trigonometric functions with: $h(x) = \text{exp}(\frac{||x||^2}{2}), l = 2, f_1 = sin, f_2 = cos$. They call it $\widehat{\text{SM}}^{trig}_m(x,y)$.
The final formula, modified using the stable random feature map mechanism proposed by Performer, is as follows:
\begin{lemma}(Positive Random Features (PRFs) for Softmax). For $x,y\in \mathbb{R}^d, z=x+y$ we have:
\begin{equation} \label{normalsm}
\begin{split}
\text{SM}(x,y)=\mathbb{E}_{\omega\~\mathcal{N}(0,\textbf{I}_d)}\left[\text{exp}(\omega^\top x-\frac{||x||^2}{2})\text{exp}(\omega^\top y-\frac{||y||^2}{2})\right]=\Lambda\mathbb{E}_{\omega\~\mathcal{N}(0,\textbf{I}_d)}\text{cosh}(\omega^\top z)
\end{split}
\end{equation}
where $\Lambda=\text{exp}(-\frac{||x||^2+||y||^2}{2})$ and \text{cosh} is hyperbolic cosine. Consequently, softmax-kernel admits a positive random feature map unbiased approximation with $h(x)=\text{exp}(-\frac{||x||^2}{2}), l=1, f_1=\text{exp}$ and $\mathcal{D}=\mathcal{N}(0, \textbf{I}_d)$ or: $h(x) = \frac{1}{\sqrt{2}}\text{exp}(-\frac{||x||^2}{2}), l=2, f_1(u)=\text{exp}(u), f_2(u)=\text{exp}(-u)$ and the same $\mathcal{D}$  (the latter for further variance reduction). We call related estimators: $\widehat{\text{SM}}^+_m$ and $\widehat{\text{SM}}^{\text{hyp+}}_m$.
\end{lemma}

Performer replaces Equation~\ref{normalsm} $\omega$ with $\sqrt{d}\frac{\omega}{||\omega||}$, they obtain the so-called \textbf{regularized softmax-kernel} SMREG which they can approximate in
a similar manner, simply changing $\mathcal{D = N} (0, \textbf{I}_d)$ to $\mathcal{D} = \text{Unif}(\sqrt{d}\mathcal{\textbf{S}}^{d-1})$, a distribution corresponding to Haar measure on the sphere of radius $\sqrt{d}$ in $\mathbb{R}^d$, obtaining estimator $\widehat{SMERG}^+_m$. 


\subsection{Notation for RetNet}
Retention introduce the mechanism that has a dual form of recurrence and
parallelism. So they can train the models in a parallel way while recurrently conducting inference.
Given input $X\in \mathbb{R}^{|x|\times d_{model}}$, they project it to one-dimensional function $v(n) = X_n · w_V$ . Consider a sequence modeling problem that maps $v(n) \rightarrow o(n)$ through states $s_n$. Let $v_n,o_n$, on denote $v(n),o(n)$ for simplicity. Retention formulate the mapping in a recurrent manner:
\begin{equation} \label{inference1}
\begin{split}
s_n = As_{n-1} + K^\top_nv_n, \ \ A\in \mathbb{R}^{d\times d}, K_n \in \mathbb{R}^{1\times d} \\ 
o_n=Q_ns_n=\sum^n_{m=1}Q_nA^{n-m}K^\top_mv_m, Q_n \in \mathbb{R}^{1\times d}
\end{split}
\end{equation}
where we map $v_n$ to the state vector $s_n$, and then implement a linear transform to encode sequence information recurrently.

Next, we make the projection $Q_n, K_n$ content-aware:
\begin{equation} \label{inference2}
\begin{split}
Q=XW_Q, K=XW_K
\end{split}
\end{equation}
where $W_{\mathcal{Q}}, W_{\mathcal{K}} \in \mathbb{R}^{d\times d}$ are learnable matrices.
We diagonalize the matrix $A=\Lambda(\gamma e^{i\theta})\Lambda^{-1}$, where $\gamma,\theta \in \mathbb{R}^d$. Then we obtain $A^{n-m}= \Theta(\gamma e^{i\theta})^{n-m}\Lambda^{-1}$.By absorbing $\Lambda$ into $W_{\mathcal{Q}}$ and $W_{\mathcal{K}}$, we can rewrite Equation~\ref{inference1} as:
\begin{equation} \label{inference3}
\begin{split}
o_n = \sum^n_{m=1}Q_n(\gamma e^{i\theta})^{n-m}K^\top_mv_m\\
=\sum^n_{m=1}(Q_n(\gamma e^{i\theta})^n)(K_m(\gamma e^i\theta)^{-m})^\top v_m
\end{split}
\end{equation}
where $Q_n(\gamma e^{i\theta})^n, K_m(\gamma e^{i\theta})^{-m}$ is known as xPos, i.e., a relative position embedding
proposed for Transformer. We further simplify $\gamma$ as a scalar, Equation~\ref{inference3} becomes:
\begin{equation} \label{inference4}
\begin{split}
o_n = \sum^n_{m=1}\gamma^{n-m}(Q_n e^{in\theta)}(K_m e^{im\theta})^{\dagger}v_m
\end{split}
\end{equation}
where $\dagger$ is the conjugate transpose. The formulation is easily parallelizable within training instances. 

In summary, Retention starts with recurrent modeling as shown in Equation~\ref{inference1}, and then derive its parallel formulation in Equation~\ref{inference3\4}. They consider the original mapping $v(n) \rightarrow o(n)$ as vectors and obtain the retention mechanism as follows.

The retention layer is defined as:
\begin{equation} \label{inference5}
\begin{split}
Q=(XW_Q)\odot \Theta, K=(XW_K)\odot \bar{\Theta}, V=XW_V \\
\Theta_n = e^{il\theta}, D_{lm}= 
\begin{cases}
\gamma^{\,n-m}, & n \ge m \\
0, & n < m
\end{cases}
, \\ \text{Retention}(X)=(QK^\top \odot D)V
\end{split}
\end{equation}
where $\bar{\Theta}$ is the complex conjugate of $\Theta$, and $D \in \mathbb{R}^{|x| \times |x|}$ combines causal masking and exponential decay along relative distance as one matrix. Similar to self-attention, the parallel representation enables us to train the models with GPUs efficiently.

the proposed mechanism can also be written as recurrent neural networks (RNNs), which is favorable for inference. For the $n$-th timestep, they recurrently obtain the output as:
\begin{equation} \label{inference6}
\begin{split}
S_n = \gamma S_{n-1}+ K^\top_nV_n\\
\text{Retention}(X_n) = Q_nS_n, \ \ n=1,...,|x|
\end{split}
\end{equation}
The $O(1)$ complexity reduces memory and inference latency while achieving equivalent results. They also discuss the inference efficiency of the compared methods. For training, RWKV’s token-mixing complexity is $O(dn)$ while Hyena’s is $O(dn log n)$ with Fast Fourier Transform acceleration. The above two methods reduce training FLOPS via employing element-wise operators to trade-off modeling capacity. In comparison with retention, the chunk-wise recurrent representation is $O(dn(b + h))$, where $b$ is the chunk size, $h$ is the head dimension, and we usually set $b = 512, h = 256$. For either large model size (i.e., larger $d$) or sequence length, the additional $b + h$ has negligible effects. So the RetNet training is quite efficient without sacrificing the modeling performance. For inference, among the compared efficient architectures, Hyena has the same complexity (i.e., $O(n)$ per step) as Transformer while the others can perform $O(1)$ decoding.

\subsection{Soundness of the Freivalds Relation}
Since $\hat{D}^{-1}$ is a diagonal matrix, the overall complexity $O(nkd)$ is determined by the $\mathbf{Q}^{'}((\mathbf{K}^{'})^\top \mathbf{V})$ computation, excluding the cost of forming $\hat{D}^{-1}$. In the following section, we demonstrate that the Freivalds’ algorithm correctly verifies the computation of $\mathbf{Q}^{'}((\mathbf{K}^{'})^\top \mathbf{V})$. Suppose the prover outputs a matrix $y'$ that has a different \textbf{structure or value} from the correct output $y=\mathbf{Q}^{'}((\mathbf{K}^{'})^\top \mathbf{V})$.
Let
\[
  s_1' = \vec{r}\ y', \qquad
  s_2 = (\vec{r}\ \mathbf{Q}^{'})((\mathbf{K}^{'})^\top \mathbf{V}).
\]
where $\mathbf{Q}^{'},\mathbf{K}^{'} \in \mathbb{R}^{n\times k}, \mathbf{V} \in \mathbb{R}^{n\times d}, \vec{r} \in \{0,1\}^n$.
% 새로운 값들로 프리발드 알고리즘이 성립할 확률.
Let $\epsilon$ denote the error probability that the Freivalds verification outputs TRUE even when it is incorrect. We claim that if $\mathbf{Q}^{'} \times ((\mathbf{K}^{'})^\top)\mathbf{V} = y'$, then $\epsilon = 0$, and if $\mathbf{Q}^{'} \times ((\mathbf{K}^{'})^\top)\mathbf{V} \neq y'$, then $\epsilon \leq \frac{1}{2}$.
\begin{itemize}
    \item \textbf{Case $\mathbf{Q}^{'} \times ((\mathbf{K}^{'})^\top)\mathbf{V} = y'$ :} The error matrix $D= \mathbf{Q}^{'} \times ((\mathbf{K}^{'})^\top)\mathbf{V} -y'$ is $0_{n\times d}$. Regardless of the choice of the vector $\vec{r}$, the equality is always satisfied because $\vec{r}D=\vec{0}$. Hence, 
    \[ \label{neqzero}
    \begin{split}
        \epsilon = \text{Pr}[\vec{r}D\neq 0] = \text{Pr}[\vec{r}0\neq \vec{0}] = 0
    \end{split}
    \]
    Therefore, the probability of error is $\epsilon=0$.
    \item \textbf{Case $\mathbf{Q}^{'} \times ((\mathbf{K}^{'})^\top)\mathbf{V} \neq y'$ : } Let $\vec{P} = \vec{r}D= (p_1,p_2,...,p_d)$, where $D = \mathbf{Q}^{'} \times ((\mathbf{K}^{'})^\top \mathbf{V}) - y'.$ Since $\mathbf{Q}^{'} \times ((\mathbf{K}^{'})^\top)\mathbf{V} \neq y',$ there exists at least one nonzero element in $D$. Let $d_{i^*,j}$ be the first nonzero entry in column $j$. For the value $\Delta=\sum_{k\neq i^*} r_kd_{kj}$. Using Bayes' Theorem, we can partition over $\Delta$.

    \[ \label{neqzero}
    \begin{split}
        p_j=\sum_{k=1}^n r_k d_{kj}= r_{i^*} d_{i^* j} + \sum_{k\neq i^*} r_k d_{kj} = r_{i^*} d_{i^* j} + \Delta \\
        \text{Pr}[p_j=0]=\text{Pr}[p_j=0|\Delta=0]\cdot \text{Pr}[\Delta=0] + \text{Pr}[p_j=0|\Delta\neq0]\cdot \text{Pr}[\Delta\neq 0] \\
        \text{Pr}[p_j=0|\Delta=0]=\text{Pr}[r_i^*=0]=\frac{1}{2} \\
    \end{split}
    \]

    For $\text{Pr}[p_j=0|\Delta\neq0]$, if $r_{i^*}d_{i^*j} + \Delta = 0$, it holds as long as $r_{i^*}=1$ (among the possible values of $r_{i^*} \in \{0,1\}$)so 
    
    \[
        \text{Pr}[p_j=0|\Delta\neq0] = \Pr[r_{i^*}=1] = \frac{1}{2}
    \]
        
    if $d_{i^*j} + \Delta \neq 0$, then $r_{i^*}$ cannot take any value to satisfy the condition, so $\text{Pr}[p_j=0|\Delta\neq0] = 0$. Therefore, $\text{Pr}[p_j=0|\Delta\neq0] \le \frac{1}{2}$.

    After taking this case into account, the probability can be rewritten as follows:

    \[ \label{neqzero2}
    \begin{split}
        \text{Pr}[p_j=0] \leq \frac{1}{2} \cdot \text{Pr}[\Delta=0] + \frac{1}{2}\cdot \text{Pr}[\Delta\neq 0] \\
        = \frac{1}{2} \cdot \text{Pr}[\Delta=0] + \frac{1}
        {2}\cdot(1-\text{Pr}[\Delta=0])\\
        =\frac{1}{2} \\
        \epsilon=\text{Pr}[\vec{P}=0| y' \neq y] \\
        = \text{Pr}[\vec{P}=0]=\text{Pr}[p_1=0 \wedge \cdots \wedge p_j=0 \wedge \cdots \wedge p_d=0] \leq \text{Pr}[p_j=0] \leq \frac{1}{2}\\
        \epsilon \leq \frac{1}{2}
    \end{split}
    \]
\end{itemize}

After $t$ independent repetitions, the $\epsilon$ becomes $(\frac{1}{2})^t$, which is negligible in $t$. Therefore, the verification protocol is sound: the probability that it outputs TRUE when $\mathbf{Q}^{'} \times ((\mathbf{K}^{'})^\top)\mathbf{V} \neq y'$ is $\epsilon \leq (\frac{1}{2})^t$.


% notation 수정 버전
\subsection{Soundness of the Freivalds Relation}
Since $\hat{D}^{-1}$ is a diagonal matrix, the overall complexity $O(nkd)$ is determined by the $\mathbf{\tilde F{A'}^+}((\mathbf{B}^{'})^\top \mathbf{V})$ computation, excluding the cost of forming $\hat{D}^{-1}$. In the following section, we demonstrate that the Freivalds’ algorithm correctly verifies the computation of $\mathbf{\tilde F{A'}^+}((\mathbf{B}^{'})^\top \mathbf{V})$. Suppose the prover outputs a matrix $y'$ that has a different \textbf{structure or value} from the correct output $y=\mathbf{\tilde F{A'}^+}((\mathbf{B}^{'})^\top \mathbf{V})$.
Let
\[
  s_1' = \vec{r}\ y', \qquad
  s_2 = (\vec{r}\ \mathbf{\tilde F{A'}^+})((\mathbf{B}^{'})^\top \mathbf{V}).
\]
where $\mathbf{\tilde F{A'}^+},\mathbf{B}^{'} \in \mathbb{R}^{n\times k}, \mathbf{V} \in \mathbb{R}^{n\times d}, \vec{r} \in \{0,1\}^n$.
% 새로운 값들로 프리발드 알고리즘이 성립할 확률.
Let $\epsilon$ denote the error probability that the Freivalds verification outputs TRUE even when it is incorrect. We claim that if $\mathbf{\tilde F{A'}^+} \times ((\mathbf{B}^{'})^\top)\mathbf{V} = y'$, then $\epsilon = 0$, and if $\mathbf{\tilde F{A'}^+} \times ((\mathbf{B}^{'})^\top)\mathbf{V} \neq y'$, then $\epsilon \leq \frac{1}{2}$.
\begin{itemize}
    \item \textbf{Case $\mathbf{\tilde F{A'}^+}((\mathbf{B}^{'})^\top \mathbf{V})=y'$ :} The error matrix $D= \mathbf{\tilde F{A'}^+}((\mathbf{B}^{'})^\top \mathbf{V}) -y'$ is $0_{n\times d}$. Regardless of the choice of the vector $\vec{r}$, the equality is always satisfied because $\vec{r}D=\vec{0}$. Hence, 
    \[ \label{neqzero}
    \begin{split}
        \epsilon = \text{Pr}[\vec{r}D\neq 0] = \text{Pr}[\vec{r}0\neq \vec{0}] = 0
    \end{split}
    \]
    Therefore, the probability of error is $\epsilon=0$.
    \item \textbf{Case $\mathbf{\tilde F{A'}^+}((\mathbf{B}^{'})^\top \mathbf{V})\neq y'$ : } Let $\vec{P} = \vec{r}D= (p_1,p_2,...,p_d)$, where $D = \mathbf{\tilde F{A'}^+}((\mathbf{B}^{'})^\top \mathbf{V}) - y'.$ Since $\mathbf{\tilde F{A'}^+}((\mathbf{B}^{'})^\top \mathbf{V})\neq y'$, there exists at least one nonzero element in $D$. Let $d_{i^*,j}$ be the first nonzero entry in column $j$. For the value $\Delta=\sum_{k\neq i^*} r_kd_{kj}$. Using Bayes' Theorem, we can partition over $\Delta$.

    \[ \label{neqzero}
    \begin{split}
        p_j=\sum_{k=1}^n r_k d_{kj}= r_{i^*} d_{i^* j} + \sum_{k\neq i^*} r_k d_{kj} = r_{i^*} d_{i^* j} + \Delta \\
        \text{Pr}[p_j=0]=\text{Pr}[p_j=0|\Delta=0]\cdot \text{Pr}[\Delta=0] + \text{Pr}[p_j=0|\Delta\neq0]\cdot \text{Pr}[\Delta\neq 0] \\
        \text{Pr}[p_j=0|\Delta=0]=\text{Pr}[r_i^*=0]=\frac{1}{2} \\
    \end{split}
    \]

    For $\text{Pr}[p_j=0|\Delta\neq0]$, if $r_{i^*}d_{i^*j} + \Delta = 0$, it holds as long as $r_{i^*}=1$ (among the possible values of $r_{i^*} \in \{0,1\}$)so 
    
    \[
        \text{Pr}[p_j=0|\Delta\neq0] = \Pr[r_{i^*}=1] = \frac{1}{2}
    \]
        
    if $d_{i^*j} + \Delta \neq 0$, then $r_{i^*}$ cannot take any value to satisfy the condition, so $\text{Pr}[p_j=0|\Delta\neq0] = 0$. Therefore, $\text{Pr}[p_j=0|\Delta\neq0] \le \frac{1}{2}$.

    After taking this case into account, the probability can be rewritten as follows:

    \[ \label{neqzero2}
    \begin{split}
        \text{Pr}[p_j=0] \leq \frac{1}{2} \cdot \text{Pr}[\Delta=0] + \frac{1}{2}\cdot \text{Pr}[\Delta\neq 0] \\
        = \frac{1}{2} \cdot \text{Pr}[\Delta=0] + \frac{1}
        {2}\cdot(1-\text{Pr}[\Delta=0])\\
        =\frac{1}{2} \\
        \epsilon=\text{Pr}[\vec{P}=0| y' \neq y] \\
        = \text{Pr}[\vec{P}=0]=\text{Pr}[p_1=0 \wedge \cdots \wedge p_j=0 \wedge \cdots \wedge p_d=0] \leq \text{Pr}[p_j=0] \leq \frac{1}{2}\\
        \epsilon \leq \frac{1}{2}
    \end{split}
    \]
\end{itemize}

After $t$ independent repetitions, the $\epsilon$ becomes $(\frac{1}{2})^t$, which is negligible in $t$. Therefore, the verification protocol is sound: the probability that it outputs TRUE when $\mathbf{\tilde F{A'}^+}((\mathbf{B}^{'})^\top \mathbf{V})\neq y'$ is $\epsilon \leq (\frac{1}{2})^t$.

\end{document}

\bibliographystyle{plain}
\bibliography{refs}